% =====================================================================
\chapter{Sistemas de numeración y coma flotante}
\label{tema:6}
% =====================================================================

Un ordenador solo almacena bits. Este tema estudia cómo se representan con ellos los enteros sin signo, los enteros con signo y los números reales, y qué se pierde por el camino. La representación posicional se fundamenta en la división euclídea de enteros; la parte de máquina sigue el capítulo~2 «Representing and Manipulating Information» de Bryant y O'Hallaron~\cite{csapp} y el artículo de Goldberg~\cite{goldberg}; los formatos de coma flotante son los de la norma IEEE~754~\cite{ieee754}.

\section{Representación posicional}

\begin{teorema}[Representación en base $b$]\label{teo:base}
Sea $b\ge2$ un entero. Todo $n\in\N$ se escribe de manera única como
\[
n=\sum_{k=0}^{m}d_kb^k=(d_md_{m-1}\dots d_1d_0)_b,\qquad d_k\in\{0,1,\dots,b-1\},\ d_m\neq0.
\]
Los $d_k$ son las \emph{cifras} de $n$ en base $b$.
\end{teorema}

La demostración es un algoritmo: las cifras de $n$ en base $b$ son los restos de dividir sucesivamente entre $b$, leídos de abajo arriba. En sentido contrario, el valor de $(d_m\dots d_0)_b$ se obtiene con el \emph{esquema de Horner}, $(\cdots((d_mb+d_{m-1})b+d_{m-2})\cdots)b+d_0$, que es la regla de Ruffini del \cref{tema:3} evaluando en $x=b$.

\addtocounter{definicion}{1}%

\begin{proposicion}[Número de cifras]\label{prop:cifras}
Un entero $n\ge1$ tiene exactamente $\lfloor\log_bn\rfloor+1$ cifras en base $b$. Con $m$ cifras en base $b$ se representan exactamente los $b^m$ enteros $0,1,\dots,b^m-1$.
\end{proposicion}

Las cifras de la parte fraccionaria de un número $x\in[0,1)$ se obtienen multiplicando sucesivamente por $b$ y tomando la parte entera: $x=\sum_{k\ge1}d_kb^{-k}=(0{,}d_1d_2\dots)_b$.

\begin{proposicion}\label{prop:fraccion-binaria}
Sea $p/q$ una fracción irreducible con $0<p<q$. Su desarrollo en base~2 es finito si y solo si $q$ es una potencia de~2.
\end{proposicion}

Por ejemplo, $0{,}1=1/10$ tiene denominador $10=2\cdot5$, de modo que su desarrollo binario es infinito: $0{,}1=(0{,}0\overline{0011})_2$. Ningún ordenador binario almacena $0{,}1$ exactamente.

\section{Binario, hexadecimal y bytes}

Un \emph{bit} es una cifra binaria y un \emph{byte} son 8 bits, que representan $2^8=256$ valores. Como $16=2^4$, cada cifra \emph{hexadecimal} ($0$--$9$, $\texttt{A}$--$\texttt{F}$) corresponde a un grupo de 4 bits, y un byte se escribe con dos cifras hexadecimales, de $\texttt{0x00}$ a $\texttt{0xFF}$. Para convertir entre binario y hexadecimal basta agrupar los bits de cuatro en cuatro desde la derecha: $(1011\,0110)_2=\texttt{0xB6}=182$.

Para cantidades de memoria se usan dos sistemas de prefijos: los decimales del SI ($1\ \text{kB}=10^3$ bytes, $1\ \text{GB}=10^9$ bytes) y los binarios de la IEC ($1\ \text{KiB}=2^{10}$ bytes, $1\ \text{GiB}=2^{30}$ bytes). La diferencia es de un $7{,}4\,\%$ en el caso de los gigas.

\section{Enteros con signo: complemento a dos}

\begin{definicion}[Complemento a dos]\label{def:c2}
Una palabra de $n$ bits $w=(a_{n-1}\dots a_0)$ representa, \emph{sin signo}, el entero $U(w)=\sum_{k=0}^{n-1}a_k2^k$ y, en \emph{complemento a dos}, el entero
\[
T(w)=-a_{n-1}2^{n-1}+\sum_{k=0}^{n-2}a_k2^k=U(w)-a_{n-1}2^n.
\]
\end{definicion}

\begin{teorema}\label{teo:c2}
Para $n\ge2$:
\begin{enumerate}
  \item $T$ es una biyección entre las palabras de $n$ bits y los enteros del intervalo $[-2^{n-1},\,2^{n-1}-1]$; $T(w)<0$ si y solo si $a_{n-1}=1$;
  \item $T(w)$ y $U(w)$ difieren en un múltiplo de $2^n$;
  \item si $\bar w$ es la palabra con todos los bits invertidos, $T(\bar w)+1=-T(w)$; en particular, el opuesto de $T(w)$ se obtiene invirtiendo los bits y sumando~1, y el resultado es representable en $n$ bits salvo si $T(w)=-2^{n-1}$.
\end{enumerate}
\end{teorema}

\addtocounter{definicion}{1}%

Por el \cref{teo:c2}(ii), la suma de enteros con signo se realiza con el mismo circuito que la suma sin signo, descartando el acarreo final: el resultado es correcto salvo que salga del rango, situación que se llama \emph{desbordamiento} (\cref{eje:6-desbordamiento}).

\section{Coma flotante}

\begin{definicion}[Número en coma flotante normalizado]\label{def:flotante}
Un formato binario de coma flotante con $e$ bits de exponente, sesgo $\beta=2^{e-1}-1$ y $p$ bits de fracción representa la palabra $(s,E,f)$, con $1\le E\le2^e-2$, como
\[
x=(-1)^s\,(1{,}f_1f_2\dots f_p)_2\cdot2^{E-\beta}.
\]
Los valores extremos de $E$ se reservan para el cero, los números subnormales, los infinitos y \texttt{NaN}~\cite{ieee754}.
\end{definicion}

\begin{table}[h]
\centering\small
\caption{Formatos binarios de coma flotante habituales. El épsilon es $\varepsilon=2^{-p}$, la distancia de~1 al siguiente número representable.}\label{tab:formatos}
\begin{tabular}{@{}lcccccc@{}}
\toprule
Formato & Bits & Signo & Exponente $e$ & Fracción $p$ & $\varepsilon=2^{-p}$ & Máximo finito \\
\midrule
fp64 (doble) & 64 & 1 & 11 & 52 & $2{,}2\cdot10^{-16}$ & $1{,}8\cdot10^{308}$ \\
fp32 (simple) & 32 & 1 & 8 & 23 & $1{,}2\cdot10^{-7}$ & $3{,}4\cdot10^{38}$ \\
fp16 (media) & 16 & 1 & 5 & 10 & $9{,}8\cdot10^{-4}$ & $65\,504$ \\
bf16 (bfloat16) & 16 & 1 & 8 & 7 & $7{,}8\cdot10^{-3}$ & $3{,}4\cdot10^{38}$ \\
\bottomrule
\end{tabular}
\end{table}

El formato bfloat16 conserva los 8 bits de exponente de fp32 y recorta la fracción a 7 bits: tiene el mismo rango que fp32 con mucha menos precisión, y sus 16 bits son exactamente los 16 bits superiores de un fp32.

\begin{proposicion}[Error de redondeo]\label{prop:redondeo}
Sea $\fl(x)$ el número del formato más cercano a $x$. Si $x$ está en el rango de los números normalizados, entonces
\[
\fl(x)=x(1+\delta)\qquad\text{con}\qquad\abs{\delta}\le u=2^{-(p+1)}=\frac\varepsilon2.
\]
\end{proposicion}

La demostración es el \cref{eje:6-redondeo}. El número $u$ se llama \emph{unidad de redondeo}: el error \emph{relativo} es pequeño, pero el absoluto crece con $\abs{x}$, porque la distancia entre números consecutivos en $[2^k,2^{k+1})$ es $2^{k-p}$~\cite{goldberg}.

\addtocounter{definicion}{1}%

\addtocounter{definicion}{1}%

